\documentclass{article}

\usepackage{amsmath,amssymb}
\usepackage{xurl}
\usepackage[hidelinks]{hyperref}
\setlength{\emergencystretch}{2em}

% Shared commands for notes in docs/paper-gaps/.

\DeclareUnicodeCharacter{2080}{\ifmmode{}_{0}\else\textsubscript{0}\fi}
\DeclareUnicodeCharacter{2081}{\ifmmode{}_{1}\else\textsubscript{1}\fi}
\DeclareUnicodeCharacter{2082}{\ifmmode{}_{2}\else\textsubscript{2}\fi}
\DeclareUnicodeCharacter{2099}{\ifmmode{}_{n}\else\textsubscript{n}\fi}

\newcommand{\Lean}{\textsc{Lean}}
\ExplSyntaxOn
\NewDocumentCommand{\leanid}{m}{
  \ifmmode
    \textsf{\detokenize{#1}}
  \else
    \group_begin:
    \urlstyle{sf}
    \tl_set:Nn \l_tmpa_tl {#1}
    \tl_replace_all:Nnn \l_tmpa_tl {\_} {_}
    \exp_args:NV \path \l_tmpa_tl
    \group_end:
  \fi
}
\ExplSyntaxOff
\providecommand{\tr}{\operatorname{tr}}
\newcommand{\tnleanIssueBase}{https://github.com/LionSR/TNLean/issues/}
\newcommand{\tnleanPrBase}{https://github.com/LionSR/TNLean/pull/}
\newcommand{\tnleanDocBase}{https://sirui-lu.com/TNLean/docs/}
\newcommand{\ghissue}[1]{\href{\tnleanIssueBase#1}{GitHub issue \##1}}
\newcommand{\ghpr}[1]{\href{\tnleanPrBase#1}{PR \##1}}
\newcommand{\doclink}[1]{\href{\tnleanDocBase#1.html}{\path{#1}}}

% Dirac notation and the identity operator, as in blueprint/src/macros/common.tex.
% \providecommand keeps note-local definitions authoritative.
\usepackage{dsfont}
\providecommand{\ket}[1]{|#1\rangle}
\providecommand{\bra}[1]{\langle #1|}
\providecommand{\braket}[2]{\langle #1 | #2 \rangle}
\providecommand{\ketbra}[2]{|#1\rangle\!\langle #2|}
\providecommand{\Id}{\mathds{1}}
\providecommand{\one}{\mathds{1}}
\providecommand{\C}{\mathbb{C}}
\providecommand{\MN}[1]{M_{#1}(\C)}
\providecommand{\MD}{M_D(\C)}

% External referencing for paper sources.  The build helper writes sanitized
% auxiliary files under build/paper-aux/ and excludes duplicated source labels.
\usepackage{xr}

% Import a generated paper auxiliary file with a caller-supplied label prefix.
% The candidate paths support builds from docs/paper-gaps/, the repository root,
% and build/paper-aux/ itself.
\newcommand{\importpaperaux}[2]{%
  \IfFileExists{../../build/paper-aux/#2.aux}{%
    \externaldocument[#1]{../../build/paper-aux/#2}%
  }{%
    \IfFileExists{build/paper-aux/#2.aux}{%
      \externaldocument[#1]{build/paper-aux/#2}%
    }{%
      \IfFileExists{#2.aux}{%
        \externaldocument[#1]{#2}%
      }{}%
    }%
  }%
}

% General external reference macros
\newcommand{\paperref}[2]{\ref{#1-#2}}
\newcommand{\papereqref}[2]{(\ref{#1-#2})}

% CPSV16 source (1606.00608 / MPDO-22-12-17-2.tex) external document and shortcuts
\importpaperaux{cpsv16-}{1606.00608/MPDO-22-12-17-2-tnlean}
\newcommand{\cpsvref}[1]{\paperref{cpsv16}{#1}}
\newcommand{\cpsveqref}[1]{\papereqref{cpsv16}{#1}}

% Verdict marker: \gapnote{<kind>}{<status>}. Kind is the policy
% classification (clarification, local-correction, scope-restriction,
% unfaithful, false-source, open-gap); status is open, wip (elimination
% underway), resolved, or historical. Machine-read by the site index;
% prints a verdict line.
\newcommand{\gapnote}[2]{\par\noindent\textbf{Verdict:} #1 (#2).\par}


\title{Global Phase Freedom in the PGVWC07 Quadratic Reconstruction System}
\author{}
\date{2026-08-28}

\begin{document}
\maketitle
\gapnote{false-source}{open}

\section{The printed assertions}

Theorem~8 of P\'erez-Garc\'ia, Verstraete, Wolf, and Cirac considers a
translation-invariant state on a finite ring with unique open-boundary
canonical form and successive reduced-state ranks bounded by $D^2$. It forms a
quadratic system from a window of $D^4$ matrices in the open-boundary canonical
form. Its second conclusion states that every solution supplies a
translation-invariant $D$-dimensional MPS representation of the state and is
related to the canonical tensor by exact unitary conjugacy
\cite[Theorem~8, lines~1154--1165]{PerezGarcia2007Matrix}.

Write $A=(A_i)_i$ for the canonical translation-invariant tensor and
$B=(B_i)_i$ for the tensor component of a solution. The printed conclusion is
therefore
\begin{align}
    V^{(N)}(B)
        &=V^{(N)}(A),
    \label{eq:pgvwc07_quadratic_reconstruction_phase_printed_state}\\
    A_i
        &=UB_iU^\dagger
        \qquad\text{for every }i
    \label{eq:pgvwc07_quadratic_reconstruction_phase_printed_conjugacy}
\end{align}
for some unitary $U$. Here $V^{(N)}(B)$ denotes the length-$N$ trace-MPS
vector generated by $B$.

The preceding uniqueness theorem makes the same exact-conjugacy assertion for
two canonical translation-invariant representations of the same finite-ring
state \cite[Theorem~7, lines~1002--1015]{PerezGarcia2007Matrix}. Both
assertions overlook the scalar phase freedom at fixed ring length.

\section{Phase symmetry of the quadratic system}

Use the zero-based notation of the blueprint and the formal quadratic-system
structure. A window has matrices $A_i^{[j]}$ for $0\leq j<m$. A solution
consists of matrices $Y_j,Z_j$ and a tensor $B$ satisfying
\begin{align}
    B_i\otimes\Id
        &=Y_jA_i^{[j]}Z_j,
    \label{eq:pgvwc07_quadratic_reconstruction_phase_system_gauge}\\
    Y_{j+1}Z_j
        &=\Id
        \qquad(0\leq j<m-1),
    \label{eq:pgvwc07_quadratic_reconstruction_phase_system_inverse}\\
    \sum_iB_iB_i^\dagger
        &=\Id.
    \label{eq:pgvwc07_quadratic_reconstruction_phase_system_normalization}
\end{align}
This is the shifted indexing of the system printed immediately before
Theorem~8 \cite[lines~1142--1152]{PerezGarcia2007Matrix}.

Let $\eta\in\C$ have unit modulus. Define
\begin{align}
    B_i^{(\eta)}
        &:=\eta B_i,
    \label{eq:pgvwc07_quadratic_reconstruction_phase_rescaled_tensor}\\
    Y_j^{(\eta)}
        &:=\overline{\eta}^{\,j}Y_j,
    \label{eq:pgvwc07_quadratic_reconstruction_phase_rescaled_left_gauge}\\
    Z_j^{(\eta)}
        &:=\eta^{j+1}Z_j.
    \label{eq:pgvwc07_quadratic_reconstruction_phase_rescaled_right_gauge}
\end{align}
Since $\overline\eta\eta=1$, the gauge equation transforms as
\begin{align}
    Y_j^{(\eta)}A_i^{[j]}Z_j^{(\eta)}
        &=\overline\eta^{\,j}\eta^{j+1}
          Y_jA_i^{[j]}Z_j
    \label{eq:pgvwc07_quadratic_reconstruction_phase_gauge_rescaling_step}\\
        &=\eta(B_i\otimes\Id)
    \label{eq:pgvwc07_quadratic_reconstruction_phase_gauge_rescaling_tensor}\\
        &=B_i^{(\eta)}\otimes\Id.
    \label{eq:pgvwc07_quadratic_reconstruction_phase_gauge_rescaling_result}
\end{align}
The internal inverse equation becomes
\begin{align}
    Y_{j+1}^{(\eta)}Z_j^{(\eta)}
        &=\overline\eta^{\,j+1}\eta^{j+1}Y_{j+1}Z_j
    \label{eq:pgvwc07_quadratic_reconstruction_phase_inverse_rescaling_step}\\
        &=\Id.
    \label{eq:pgvwc07_quadratic_reconstruction_phase_inverse_rescaling_result}
\end{align}
Finally, the normalization is unchanged:
\begin{align}
    \sum_iB_i^{(\eta)}(B_i^{(\eta)})^\dagger
        &=\eta\overline\eta\sum_iB_iB_i^\dagger
    \label{eq:pgvwc07_quadratic_reconstruction_phase_normalization_rescaling_step}\\
        &=\Id.
    \label{eq:pgvwc07_quadratic_reconstruction_phase_normalization_rescaling_result}
\end{align}
Thus every solution belongs to a full unit-modulus phase family of solutions.
This freedom is present in the quadratic system itself and is not a virtual
similarity gauge.

\section{Effect on a finite-ring state}

Every length-$N$ coefficient is homogeneous of degree $N$ in the local tensor.
Therefore
\begin{align}
    V^{(N)}(B^{(\eta)})
        &=\eta^N V^{(N)}(B).
    \label{eq:pgvwc07_quadratic_reconstruction_phase_finite_ring_effect}
\end{align}
Suppose $V^{(N)}(B)$ is the prescribed nonzero state. If $\eta^N\neq1$, then
(\ref{eq:pgvwc07_quadratic_reconstruction_phase_finite_ring_effect}) shows that
$B^{(\eta)}$ is another solution but does not represent that state. This
contradicts the unrestricted representation claim
(\ref{eq:pgvwc07_quadratic_reconstruction_phase_printed_state}).

If $\eta^N=1$, then $B^{(\eta)}$ does represent the same finite-ring state.
The exact-conjugacy claim can nevertheless fail. Take $d=D=1$, let the single
canonical matrix be
\begin{align}
    A_0
        &=1,
    \label{eq:pgvwc07_quadratic_reconstruction_phase_product_tensor}
\end{align}
and choose $N>3$. This tensor represents a nonzero product state, satisfies
condition C1 with onset $L_0=1$, has rank-one reduced states, and has no
Schmidt-value degeneracy. Let $\eta\neq1$ be an $N$th root of unity and set
\begin{align}
    B_0
        &=\eta A_0.
    \label{eq:pgvwc07_quadratic_reconstruction_phase_product_rephasing}
\end{align}
The phase construction above makes $B$ the tensor component of a solution, and
(\ref{eq:pgvwc07_quadratic_reconstruction_phase_finite_ring_effect}) gives
$V^{(N)}(B)=V^{(N)}(A)$. In bond dimension one, unitary conjugation acts
trivially:
\begin{align}
    UB_0U^\dagger
        &=B_0
    \label{eq:pgvwc07_quadratic_reconstruction_phase_scalar_conjugation}\\
        &\neq A_0.
    \label{eq:pgvwc07_quadratic_reconstruction_phase_scalar_conjugation_failure}
\end{align}
Hence exact unitary conjugacy fails even after restricting to solutions that
represent the same fixed-length state. This also gives a finite-ring phase
counterexample to the exact conclusion printed in Theorem~7.

\section{Scope of the endpoint-cut argument}

The uniqueness argument compares open-boundary representations over the
$D^4$-site window. The complete dimension argument requires one intertwiner at
each of the $D^4+1$ cuts and one relation across each of the $D^4$ sites. The
off-by-one correction is recorded separately in
\path{docs/paper-gaps/pgvwc07_intertwiner_chain_off_by_one.tex}.

The current quadratic-system structure produces only the internal-cut
intertwiners. The two endpoint intertwiners must come from uniqueness of the
open-boundary representation. That comparison is available only after the
solution tensor is already known to represent the same state as the canonical
tensor. It cannot prove that an arbitrary algebraic solution represents the
state, even up to phase.

Once the same-state hypothesis supplies the endpoint cuts, the corrected chain
argument gives a nonzero matrix $R$ and a nonzero scalar $x$ satisfying
\begin{align}
    RB_i
        &=x^{-1}A_iR
        \qquad\text{for every }i.
    \label{eq:pgvwc07_quadratic_reconstruction_phase_intertwiner}
\end{align}
The normalization and fixed-point argument then gives
\begin{align}
    |x|
        &=1,
    \label{eq:pgvwc07_quadratic_reconstruction_phase_intertwiner_modulus}\\
    RR^\dagger
        &=c\Id
        \qquad\text{for some }c>0.
    \label{eq:pgvwc07_quadratic_reconstruction_phase_intertwiner_norm}
\end{align}
Thus $U=c^{-1/2}R$ is unitary and
\begin{align}
    A_i
        &=xUB_iU^\dagger
        \qquad\text{for every }i.
    \label{eq:pgvwc07_quadratic_reconstruction_phase_unit_phase_conjugacy}
\end{align}
Because $A$ and $B$ generate the same nonzero length-$N$ state,
(\ref{eq:pgvwc07_quadratic_reconstruction_phase_unit_phase_conjugacy}) implies
\begin{align}
    V^{(N)}(A)
        &=x^NV^{(N)}(B)
    \label{eq:pgvwc07_quadratic_reconstruction_phase_same_state_phase_relation}\\
        &=V^{(N)}(B),
    \label{eq:pgvwc07_quadratic_reconstruction_phase_same_state_equality}\\
    x^N
        &=1.
    \label{eq:pgvwc07_quadratic_reconstruction_phase_root_condition}
\end{align}
The last conclusion uses that the state vector is nonzero.

\section{Corrected theorem}

The mathematically justified reconstruction statement keeps the existence
clause of Theorem~8 and restricts its converse. Let $B$ be the tensor of a
solution whose length-$N$ trace-MPS representation is the prescribed state.
Then there are a unitary $U$ and a scalar $\xi\in\C$ such that
\begin{align}
    |\xi|
        &=1,
    \label{eq:pgvwc07_quadratic_reconstruction_phase_corrected_unit_modulus}\\
    \xi^N
        &=1,
    \label{eq:pgvwc07_quadratic_reconstruction_phase_corrected_root_condition}\\
    A_i
        &=\xi UB_iU^\dagger
        \qquad\text{for every }i.
    \label{eq:pgvwc07_quadratic_reconstruction_phase_corrected_conjugacy}
\end{align}
Equivalently, define the explicitly rephased tensor by
\begin{align}
    \widehat B_i
        &:=\xi B_i.
    \label{eq:pgvwc07_quadratic_reconstruction_phase_corrected_rephasing}
\end{align}
Then
\begin{align}
    V^{(N)}(\widehat B)
        &=V^{(N)}(B),
    \label{eq:pgvwc07_quadratic_reconstruction_phase_rephased_same_state}\\
    A_i
        &=U\widehat B_iU^\dagger
        \qquad\text{for every }i.
    \label{eq:pgvwc07_quadratic_reconstruction_phase_rephased_exact_conjugacy}
\end{align}
The first relation is unit gauge-phase equivalence between $B$ and $A$, while
the second is exact gauge equivalence after the displayed scalar
rephasing.\footnote{The corresponding formal predicates are
    \leanid{MPSTensor.UnitGaugePhaseEquiv} and
    \leanid{MPSTensor.GaugeEquiv} in \doclink{TNLean/MPS/Defs}.}

\section{Open classification gap and verdict}

The phase family proves that the unrestricted conclusion of Theorem~8 and the
exact finite-ring conclusion of Theorem~7 are false as printed. The corrected
same-state theorem above is the conclusion supported by the endpoint-cut
argument.

A stronger classification remains open: the current quadratic equations have
not been shown to imply that every solution is phase-gauge equivalent to the
canonical tensor. The known phase-rescaled solutions do not show that they are
the only spurious solutions. Proving such a result would require either a new
argument deriving endpoint compatibility from the internal equations or an
enlarged reconstruction system that includes endpoint data explicitly.

The blueprint therefore retains the algebraic reconstruction route but does
not claim that an arbitrary solution represents the prescribed state. Its
converse is restricted to same-state solutions and records the unavoidable
$N$th-root phase before exact rephasing.

\bibliographystyle{alpha}
\bibliography{references}

\end{document}